\documentclass[10pt,a4paper]{amsart}
\usepackage{amsmath,amssymb,amsthm,graphicx}
% For commutative diagrams we use Xy-pic
\usepackage[all]{xy}

% We use 2cell for 2-commutative diagrams.
\xyoption{2cell}
% Theorem environments.
%
\theoremstyle{plain}
\newtheorem{theorem}[subsection]{Theorem}
\newtheorem{proposition}[subsection]{Proposition}
\newtheorem{lemma}[subsection]{Lemma}

\theoremstyle{definition}
\newtheorem{definition}[subsection]{Definition}
\newtheorem{example}[subsection]{Example}
\newtheorem{exercise}[subsection]{Exercise}
\newtheorem{situation}[subsection]{Situation}

\theoremstyle{remark}
\newtheorem{remark}[subsection]{Remark}
\newtheorem{remarks}[subsection]{Remarks}

\numberwithin{equation}{subsection}

% Macros
%
\def\lim{\mathop{\mathrm{lim}}\nolimits}
\def\colim{\mathop{\mathrm{colim}}\nolimits}
\def\Spec{\mathop{\mathrm{Spec}}}
\def\Hom{\mathop{\mathrm{Hom}}\nolimits}
\def\Ext{\mathop{\mathrm{Ext}}\nolimits}
\def\SheafHom{\mathop{\mathcal{H}\!\mathit{om}}\nolimits}
\def\SheafExt{\mathop{\mathcal{E}\!\mathit{xt}}\nolimits}
\def\Sch{\mathit{Sch}}
\def\Mor{\mathop{\mathrm{Mor}}\nolimits}
\def\Ob{\mathop{\mathrm{Ob}}\nolimits}
\def\Sh{\mathop{\mathit{Sh}}\nolimits}
\def\NL{\mathop{N\!L}\nolimits}
\def\CH{\mathop{\mathrm{CH}}\nolimits}
\def\proetale{{pro\text{-}\acute{e}tale}}
\def\etale{{\acute{e}tale}}
\def\QCoh{\mathit{QCoh}}
\def\Ker{\mathop{\mathrm{Ker}}}
\def\Im{\mathop{\mathrm{Im}}}
\def\Coker{\mathop{\mathrm{Coker}}}
\def\Coim{\mathop{\mathrm{Coim}}}

% Boxtimes
%
\DeclareMathSymbol{\boxtimes}{\mathbin}{AMSa}{"02}

%
% Macros for moduli stacks/spaces
%
\def\QCohstack{\mathcal{QC}\!\mathit{oh}}
\def\Cohstack{\mathcal{C}\!\mathit{oh}}
\def\Spacesstack{\mathcal{S}\!\mathit{paces}}
\def\Quotfunctor{\mathrm{Quot}}
\def\Hilbfunctor{\mathrm{Hilb}}
\def\Curvesstack{\mathcal{C}\!\mathit{urves}}
\def\Polarizedstack{\mathcal{P}\!\mathit{olarized}}
\def\Complexesstack{\mathcal{C}\!\mathit{omplexes}}
% \Pic is the operator that assigns to X its picard group, usage \Pic(X)
% \Picardstack_{X/B} denotes the Picard stack of X over B
% \Picardfunctor_{X/B} denotes the Picard functor of X over B
\def\Pic{\mathop{\mathrm{Pic}}\nolimits}
\def\Picardstack{\mathcal{P}\!\mathit{ic}}
\def\Picardfunctor{\mathrm{Pic}}
\def\Deformationcategory{\mathcal{D}\!\mathit{ef}}
\title{Cohomology of a profinite group in positive degrees is torsion and vanishes on Q-vector spaces}
\author{The Stacks Project Authors}
\date{stacks-project:etale-cohomology.tex:local-snapshot-20261011; GNU FDL 1.2 or later}
\begin{document}
\maketitle
\noindent\textit{Editorial note.} All delivered mathematical text is the author's own, copied from the pinned Stacks Project chapter file etale-cohomology.tex; no statement or proof step was rewritten, reordered or removed.\par\medskip
\noindent\textit{Editorial note.} The proof is the author's complete argument for both parts: part (1) is reduced by dimension shifting and then computed, and part (2) is the author's own one-line consequence of part (1).\par\medskip
\noindent\textit{Editorial note.} The label lemma-profinite-group-cohomology-torsion occurs in no other staged Stacks chapter, so no cross-chapter restatement had to be weighed for this card.\par\medskip
\noindent\textit{Editorial note.} The preamble is the chapter's own native preamble excerpt already used for this source: the Xy-pic setup (native preamble lines 14-18) and the native theorem declarations with the mathematical macros (native preamble lines 154-217).\par\medskip
\noindent\textit{Allowed standard prerequisites (editorial).}\begin{itemize}
\item Continuous cohomology of a profinite group and its long exact sequence. Same chapter, the definition of H\textasciicircum{}i(G, M) for a profinite group G and the associated long exact cohomology sequence. Supplies the functors H\textasciicircum{}i(G, -) and the exact sequence from which the element of H\textasciicircum{}1(G, M) is read off.
\item Dimension shifting. Same chapter, the dimension shifting argument for the cohomology of a profinite group. Reduces statement (1) in degrees i > 0 to the case i = 1, which is the only degree the counted proof computes.
\item Injective objects of the category of G-modules. Same chapter, the existence of enough injectives in the category of G-modules. Supplies the embedding 0 -> M -> I -> N -> 0 with I injective that the proof uses to represent classes in H\textasciicircum{}1(G, M).
\item Open stabilisers in a profinite group. Same chapter, the stabiliser of an element of a discrete G-module under the action of a profinite group is open and hence of finite index. Turns the stabiliser U of the chosen lift x into a finite-index subgroup, producing the integer r that kills the cohomology class.
\end{itemize}\medskip
\begin{lemma}
\label{lemma-profinite-group-cohomology-torsion}
Let $G$ be a profinite topological group.
Then
\begin{enumerate}
\item $H^i(G, M)$ is torsion for $i > 0$ and any $G$-module $M$, and
\item $H^i(G, M) = 0$ if $M$ is a $\mathbf{Q}$-vector space.
\end{enumerate}
\end{lemma}

\begin{proof}
Proof of (1). By dimension shifting we see that it suffices
to show that $H^1(G, M)$ is torsion for every $G$-module $M$.
Choose an exact sequence $0 \to M \to I \to N \to 0$ with $I$
an injective object of the category of $G$-modules.
Then any element of $H^1(G, M)$ is the image of an element
$y \in N^G$. Choose $x \in I$ mapping to $y$.
The stabilizer $U \subset G$ of $x$ is open, hence
has finite index $r$. Let $g_1, \ldots, g_r \in G$ be a system
of representatives for $G/U$. Then $\sum g_i(x)$ is an invariant
element of $I$ which maps to $ry$. Thus $r$ kills the element
of $H^1(G, M)$ we started with. Part (2) follows as then
$H^i(G, M)$ is both a $\mathbf{Q}$-vector space and torsion.
\end{proof}

\end{document}
